% !TEX program = xelatex
% AI-effects 프로젝트 --- 한국 직업(SOC)별 AI 대화량과 observed exposure의 직업 매칭 범위.
% 표 본문: results/table/20_{units,match,match_months,bins,by_major,reverse,kor_not_in_exposure}.tex
%          (code/04_analysis_kor_soc_exposure_coverage.do가 생성)
% 관련: report 19, code/03_merge_sensortower_aei_kor_soc_conversations.do
% 컴파일: xelatex -> biber -> xelatex -> xelatex (kotex, biblatex/biber 필요)
\documentclass[11pt]{article}

\usepackage{fontspec}
\usepackage{kotex}
\usepackage[a4paper,margin=2.2cm]{geometry}
\usepackage{longtable}
\usepackage{booktabs}
\usepackage{array}
\usepackage{amsmath}
\usepackage{amssymb}
\usepackage{xcolor}
\usepackage{hyperref}
\usepackage{enumitem}
\usepackage{titlesec}
\usepackage[backend=biber,style=authoryear,maxcitenames=2,uniquelist=false,sorting=nyt]{biblatex}
\addbibresource{../../reference/references.bib}

\hypersetup{colorlinks=true,linkcolor=blue!50!black,citecolor=blue!50!black,urlcolor=blue!50!black}
\newcolumntype{L}[1]{>{\raggedright\arraybackslash}p{#1}}
\newcolumntype{R}[1]{>{\raggedleft\arraybackslash}p{#1}}
\setlength{\LTleft}{0pt}
\setlength{\LTright}{0pt}
\renewcommand{\arraystretch}{1.15}
\setlength{\tabcolsep}{4pt}

\title{한국 AI 대화량과 observed exposure의 직업 매칭 범위\\[2mm]
\large --- SOC 분류 단위, 매칭되지 않는 직업의 원인, 분석상 유의점 ---}
\author{AI-effects 프로젝트}
\date{\today}

\begin{document}
\maketitle

\begin{abstract}
\noindent 본 문서는 report 19의 한국 직업별 업무용 AI 대화량을 직업별 노출도 지표인 observed exposure\parencite{massenkoff-2026-labor-market-impacts-of-ai}와 결합할 때 직업이 얼마나 매칭되는지 정리한다. 한국 대화량의 직업 비중은 AEI release 2026-06-26\parencite{massenkoff-2026-cadences}에서 가져온다. 이 자료의 세부 직업은 SOC 세부 직업(6자리)보다 한 단계 더 세분된 O*NET-SOC 8자리 코드다. 2026년 5월 한국에는 424개 코드가 공개되며, 6자리로 묶으면 365개 SOC 세부 직업이 된다. 같은 달 global에는 718개 코드, 614개 SOC 세부 직업이 공개된다. observed exposure 파일(\texttt{job\_exposure.csv})은 SOC 2018 세부 직업 756개 단위다. 이 756개 가운데 한국 대화량과 매칭되는 직업은 342개(45.2\%)에 그친다. 나머지는 global에만 공개된 직업 236개와 global에서도 공개되지 않은 직업 178개다. 매칭되지 않는 주된 이유는 AEI가 대화가 적은 셀을 공개하지 않는 데 있다. 이 기준은 표본이 작은 한국에서 더 많은 셀을 걸러낸다. 또한 Claude 대화는 사무$\cdot$전문직 과업에 몰려 있고, 건설$\cdot$생산$\cdot$운송 직업에는 거의 없다. observed exposure도 같은 Claude 사용 자료로 만들어지므로 같은 직업들이 0으로 나타난다. 매칭되지 않은 414개 직업 가운데 323개는 exposure가 0이다. 직업 수로는 절반이 넘게 빠지지만, 이 414개 직업이 차지하는 한국 대화 비중은 공개되지 않은 나머지 1.92\% 안에 있다. 반대 방향의 손실이 오히려 더 크다. 한국에 공개된 직업 중 23개(한국 대화의 3.21\%)는 \texttt{job\_exposure.csv}에 없다.
\end{abstract}

\tableofcontents

% ============================================================
\section{목적과 질문}
% ============================================================

report 19는 한국의 월간 업무용 AI 대화량을 AEI의 한국 직업 비중으로 나누었다. 이 직업별 대화량(usage)을 노출도(exposure) 지표와 직업 단위로 결합하는 것이 다음 분석 단계다. 본 문서는 결합에 앞서 다음 세 질문에 답한다.
\begin{enumerate}[leftmargin=*]
  \item 한국 대화량과 observed exposure는 각각 SOC의 어느 단위(대분류$\cdot$소분류$\cdot$broad$\cdot$세부 직업)로 분류되어 있는가?
  \item observed exposure의 756개 직업 가운데 한국 대화량과 매칭되는 직업은 몇 개이며, 왜 절반이 넘게 매칭되지 않는가?
  \item 매칭되지 않는 직업이 대화량 기준으로 얼마나 중요하며, 분석에서 어떻게 다루어야 하는가?
\end{enumerate}

% ============================================================
\section{원자료}
% ============================================================

\begin{longtable}{L{3.4cm}L{5.0cm}L{7.4cm}}
\caption{원자료}\label{tab:raw}\\
\toprule
자료 & 파일 & 내용 \\
\midrule
\endfirsthead
\toprule
자료 & 파일 & 내용 \\
\midrule
\endhead
AEI release 2026-06-26 & \texttt{aei\_claude\_ai\_2026-06-26.csv} & \texttt{category\_name} = \texttt{soc\_occupation}, \texttt{hierarchy\_level} = 0, \texttt{metric\_id} = \texttt{pct}. \texttt{geo\_id} = KOR(country), GLOBAL(global). 2026-04, 2026-05 \\
AEI labor market impacts & \texttt{job\_exposure.csv} & \texttt{occ\_code}, \texttt{title}, \texttt{observed\_exposure}(직업별 1개 값, 시점 구분 없음) \\
\bottomrule
\end{longtable}

\subsection{AEI 직업(SOC) 비중}

AEI의 \texttt{soc\_occupation} 범주는 대화 내용이 해당하는 O*NET 과업의 직업을 기준으로 대화 비중(\texttt{pct})을 공개한다. 즉 사용자의 직업이 아니라 \emph{과업의 직업}이다. 위계는 두 단계뿐이다. \texttt{hierarchy\_level} = 1은 SOC 대분류(Major group) 22개이고, \texttt{hierarchy\_level} = 0은 세부 직업(Detailed Occupation)이다. SOC의 소분류(Minor group)와 broad occupation 단위는 공개되지 않는다.

자료 설명서는 공개 기준을 다음과 같이 밝힌다. 셀은 집계 기준(aggregation threshold)과 지역 표본 하한(geography sample floor)을 모두 넘을 때만 공개된다. 지역 표본 하한은 해당 지역이 그 기간에 확보해야 하는 최소 표본 대화 수다. 따라서 행이 없다는 것은 셀이 공개되지 않았다는 뜻이지, 값이 0이라는 뜻은 아니다. 국가 단위에서는 세부 직업에 대해 \texttt{pct}만 공개되고, global에서는 세부 직업에도 모든 지표가 공개된다.

\subsection{observed exposure}

\textcite{massenkoff-2026-labor-market-impacts-of-ai}의 observed exposure는 다음 세 자료를 결합한 직업 단위 지표다.
\begin{itemize}[leftmargin=*]
  \item \textbf{이론적 노출}: \textcite{eloundou-2023-gpts-are-gpts-an-early}의 과업별 등급 $\beta$.
  \item \textbf{실제 사용}: AEI의 Claude 사용 자료. 본문은 2025년 8월$\cdot$11월 자료를 쓴다.
  \item \textbf{과업 시간 비중}: O*NET 과업에 쓰는 시간의 비중.
\end{itemize}
과업 $t$의 업무 사용량 $\text{WorkUsage}_t$는 Claude.ai의 업무용 대화와 1P API 트래픽의 합이다. 과업 단위 노출은
\[
\tilde r_t = \mathbf{1}\{\text{WorkUsage}_t \ge 100\}\times \mathbf{1}\{\beta_t \ge 0.5\}\times \alpha_t,\qquad \alpha_t = \tfrac12 + \tfrac12\,\text{AutomationShare}_t
\]
이고, 직업 $o$의 값은 과업 시간 비중 $w_t$로 가중한 평균 $R_o=\sum_t w_t \tilde r_t / \sum_t w_t$이다. 업무 사용량이 100건(전체 트래픽의 0.0025\%)에 못 미치는 과업은 노출이 0으로 처리된다. 따라서 observed exposure가 0인 직업은 대부분 ``Claude 업무 사용이 거의 관측되지 않은 직업''이다. 이 점은 4절에서 매칭 실패의 원인을 해석하는 데 중요하다.

% ============================================================
\section{직업 분류 단위}
% ============================================================

\begin{longtable}{L{6.0cm}rrrrr}
\caption{직업 분류 단위와 공개 범위}\label{tab:units}\\
\toprule
 & 한국 & 한국 & global & global & \texttt{job\_} \\
항목 & 2026-04 & 2026-05 & 2026-04 & 2026-05 & \texttt{exposure} \\
\midrule
\endfirsthead
\toprule
항목 & 한국 04 & 한국 05 & global 04 & global 05 & exposure \\
\midrule
\endhead
% 04_analysis_kor_soc_exposure_coverage.do가 생성. 직접 고치지 말 것.
O*NET-SOC 8자리 코드 수 & 389 & 424 & 696 & 718 & -- \\
\quad 끝자리 .00(= SOC 세부 직업) & 315 & 345 & 573 & 592 & -- \\
\quad 끝자리 .01 이상(O*NET 세분) & 74 & 79 & 123 & 126 & -- \\
SOC 세부 직업(6자리) 수 & 336 & 365 & 595 & 614 & 756 \\
\quad .00 없이 .01 이상으로만 나타남 & 21 & 20 & 22 & 22 & -- \\
SOC 대분류 수 & 21 & 22 & 22 & 22 & 22 \\
공개 세부 직업 pct 합계(\%) & 96.16 & 98.08 & 94.37 & 98.51 & -- \\

\bottomrule
\end{longtable}

\paragraph{AEI 세부 직업은 O*NET-SOC 8자리 코드다.} AEI의 세부 직업 코드는 \texttt{15-1252.00}처럼 SOC 6자리 뒤에 두 자리가 붙은 O*NET-SOC 코드다. 끝자리가 \texttt{.00}인 코드는 SOC 세부 직업과 같은 단위다. \texttt{.01} 이상은 O*NET이 하나의 SOC 세부 직업을 더 세분한 직업이다(예: \texttt{11-1011.03} Chief Sustainability Officers). 2026년 5월 한국의 424개 코드는 \texttt{.00} 345개와 \texttt{.01} 이상 79개로 이루어진다. 앞 6자리로 묶으면 365개 SOC 세부 직업이 된다. 이 가운데 20개는 \texttt{.00} 코드 없이 \texttt{.01} 이상 코드로만 나타난다. 같은 달 global은 718개 코드, 614개 SOC 세부 직업이다. report 19의 \texttt{03\_merge\_sensortower\_aei\_kor\_soc\_conversations.do}는 대화량을 8자리 코드 단위로 배분했다. 따라서 SOC 4단계 위계 중 어느 것과도 정확히 같지 않고, 세부 직업보다 한 단계 아래다.

\paragraph{observed exposure는 SOC 2018 세부 직업(6자리)이다.} \texttt{job\_exposure.csv}의 756개 \texttt{occ\_code}는 모두 \texttt{XX-XXXX} 형식이며 중복이 없다. 끝자리가 0인 broad occupation 코드는 없고, 22개 대분류가 모두 있다. \texttt{15-1252} Software Developers와 \texttt{15-2051} Data Scientists가 있고 SOC 2010 코드(\texttt{15-1132} 등)는 없으므로 SOC 2018 기준이다. SOC 2018의 세부 직업 867개 가운데 일부다.

\paragraph{매칭 키.} 두 자료는 AEI 코드의 앞 7자리(\texttt{XX-XXXX})와 \texttt{occ\_code}로 연결한다. 한 SOC 세부 직업에 8자리 코드가 여럿 있으면 \texttt{pct}를 합한다. 이하 모든 집계는 이 6자리 단위다.

% ============================================================
\section{매칭 결과}
% ============================================================

\texttt{job\_exposure.csv}의 756개 직업을 2026년 5월 AEI 공개 여부에 따라 세 집단으로 나눈다. 한국에 공개된 직업은 모두 global에도 공개되어 있으므로 세 집단은 서로 겹치지 않는다(do 파일에서 확인).
\begin{enumerate}[leftmargin=*,label=(\arabic*)]
  \item \textbf{한국 매칭}: 한국에 \texttt{pct}가 공개된 직업.
  \item \textbf{global에만 공개}: global에는 공개되었으나 한국에는 공개되지 않은 직업.
  \item \textbf{global에서도 미공개}: global에서도 공개되지 않은 직업.
\end{enumerate}

\begin{longtable}{L{4.2cm}rrrrrrrr}
\caption{\texttt{job\_exposure} 756개 직업의 매칭 결과(2026년 5월)}\label{tab:match}\\
\toprule
 & 직업 & 비율 & \multicolumn{2}{c}{observed exposure} & \multicolumn{2}{c}{exposure = 0} & 한국 & global \\
집단 & 수 & (\%) & 평균 & 중앙값 & 수 & (\%) & pct(\%) & pct(\%) \\
\midrule
\endfirsthead
\toprule
집단 & 수 & \% & 평균 & 중앙값 & 0 수 & 0 \% & 한국 pct & global pct \\
\midrule
\endhead
% 04_analysis_kor_soc_exposure_coverage.do가 생성. 직접 고치지 말 것.
(1) 한국 매칭 & 342 & 45.2 & 0.1440 & 0.0815 & 88 & 25.7 & 94.87 & 93.45 \\
(2) global에만 공개(한국 미공개) & 236 & 31.2 & 0.0362 & 0.0000 & 154 & 65.3 & -- & 0.83 \\
(3) global에서도 미공개 & 178 & 23.5 & 0.0023 & 0.0000 & 169 & 94.9 & -- & -- \\
\midrule
합계 & 756 & 100.0 & 0.0770 & 0.0000 & 411 & 54.4 & 94.87 & 94.28 \\

\bottomrule
\end{longtable}
\noindent 주: 한국$\cdot$global pct는 해당 집단 직업의 AEI 2026년 5월 \texttt{pct} 합계. ``--''는 공개된 값이 없음.

한국과 매칭되는 직업은 342개(45.2\%)다. 나머지 414개(54.8\%)는 global에만 공개된 236개와 global에서도 공개되지 않은 178개다. 세 집단의 observed exposure는 뚜렷하게 다르다. 평균이 각각 0.144, 0.036, 0.002이고, exposure가 0인 직업의 비율은 25.7\%, 65.3\%, 94.9\%다.

% ============================================================
\section{매칭되지 않는 원인}
% ============================================================

\subsection{AEI의 공개 기준과 한국 표본 크기}

첫째 원인은 AEI의 공개 기준이다. 세부 직업 셀은 집계 기준과 지역 표본 하한을 넘을 때만 공개된다. 한국의 Claude.ai 표본은 global보다 훨씬 작다. 그래서 global에서는 기준을 넘는 직업 236개가 한국에서는 기준에 못 미쳐 빠진다. 이 236개 직업이 global 대화에서 차지하는 비중은 합쳐서 0.83\%다(표~\ref{tab:match}). 개별 직업으로는 대부분 global에서도 작은 직업이다.

매칭 수가 표본에 따라 움직인다는 점은 월별 비교에서도 확인된다. 한국 공개 세부 직업은 2026년 4월 389개(6자리 336개)에서 5월 424개(365개)로 늘었다. 이에 따라 매칭 직업도 317개에서 342개로 25개 늘었다(표~\ref{tab:months}). 4월에는 농림어업(45) 대분류에서 공개된 세부 직업이 하나도 없었다(표~\ref{tab:units}의 대분류 수 21).

\begin{longtable}{L{6.4cm}rrr}
\caption{집단별 직업 수: 2026년 4월 대 5월}\label{tab:months}\\
\toprule
집단 & 2026-04 & 2026-05 & 변화 \\
\midrule
\endfirsthead
\toprule
집단 & 2026-04 & 2026-05 & 변화 \\
\midrule
\endhead
% 04_analysis_kor_soc_exposure_coverage.do가 생성. 직접 고치지 말 것.
(1) 한국 매칭 & 317 & 342 & 25 \\
(2) global에만 공개(한국 미공개) & 242 & 236 & -6 \\
(3) global에서도 미공개 & 197 & 178 & -19 \\

\bottomrule
\end{longtable}

\subsection{Claude 대화의 직업 집중}

둘째 원인은 Claude 대화가 일부 직업에 몰려 있다는 점이다. global에서도 178개 직업이 공개되지 않는다. 이 직업들은 전 세계 Claude.ai 대화에서도 해당 과업으로 분류되는 대화가 공개 기준에 못 미칠 만큼 적다. 표~\ref{tab:major}에서 보듯 매칭 실패는 특정 대분류에 몰려 있다. 건설$\cdot$채굴(56개 중 한국 매칭 1개), 생산(98개 중 9개), 운송$\cdot$자재 이동(44개 중 4개), 농림어업(11개 중 1개), 설치$\cdot$정비$\cdot$수리(50개 중 8개)가 그렇다. 반면 법률(7개 중 7개), 교육$\cdot$도서관(57개 중 55개), 컴퓨터$\cdot$수학(21개 중 20개), 경영(35개 중 30개)은 거의 모두 매칭된다. 직업 수가 많은 생산$\cdot$건설$\cdot$운송$\cdot$정비 대분류(합계 248개)의 매칭 실패만으로 매칭되지 않은 414개의 절반을 넘는다(226개).

\begin{longtable}{L{5.6cm}rrrrrrr}
\caption{대분류별 매칭 결과(2026년 5월)}\label{tab:major}\\
\toprule
 & 직업 & & & & 한국 매칭 & \multicolumn{2}{c}{평균 exposure} \\
대분류 & 수 & (1) & (2) & (3) & 비율(\%) & (1) & (2)+(3) \\
\midrule
\endfirsthead
\toprule
대분류 & 수 & (1) & (2) & (3) & 매칭(\%) & exp. (1) & exp. (2)+(3) \\
\midrule
\endhead
% 04_analysis_kor_soc_exposure_coverage.do가 생성. 직접 고치지 말 것.
11 경영 & 35 & 30 & 5 & 0 & 85.7 & 0.104 & 0.057 \\
13 사업·금융 운영 & 29 & 24 & 4 & 1 & 82.8 & 0.211 & 0.013 \\
15 컴퓨터·수학 & 21 & 20 & 1 & 0 & 95.2 & 0.373 & 0.486 \\
17 건축·공학 & 34 & 27 & 6 & 1 & 79.4 & 0.048 & 0.000 \\
19 생명·물리·사회과학 & 46 & 28 & 16 & 2 & 60.9 & 0.167 & 0.052 \\
21 지역사회·사회서비스 & 12 & 7 & 5 & 0 & 58.3 & 0.056 & 0.046 \\
23 법률 & 7 & 7 & 0 & 0 & 100.0 & 0.216 & -- \\
25 교육·도서관 & 57 & 55 & 2 & 0 & 96.5 & 0.134 & 0.013 \\
27 예술·디자인·엔터테인먼트·스포츠·미디어 & 34 & 24 & 9 & 1 & 70.6 & 0.181 & 0.050 \\
29 보건의료 전문·기술 & 65 & 29 & 28 & 8 & 44.6 & 0.044 & 0.039 \\
31 보건의료 지원 & 16 & 4 & 8 & 4 & 25.0 & 0.171 & 0.009 \\
33 보안 서비스 & 21 & 3 & 11 & 7 & 14.3 & 0.041 & 0.007 \\
35 음식 조리·서빙 & 15 & 3 & 12 & 0 & 20.0 & 0.024 & 0.006 \\
37 건물·부지 청소·관리 & 8 & 1 & 3 & 4 & 12.5 & 0.044 & 0.022 \\
39 개인 돌봄·서비스 & 27 & 13 & 11 & 3 & 48.1 & 0.012 & 0.026 \\
41 판매 & 20 & 15 & 5 & 0 & 75.0 & 0.249 & 0.197 \\
43 사무·행정 지원 & 50 & 29 & 18 & 3 & 58.0 & 0.264 & 0.089 \\
45 농림어업 & 11 & 1 & 4 & 6 & 9.1 & 0.000 & 0.002 \\
47 건설·채굴 & 56 & 1 & 15 & 40 & 1.8 & 0.000 & 0.004 \\
49 설치·정비·수리 & 50 & 8 & 23 & 19 & 16.0 & 0.011 & 0.006 \\
51 생산 & 98 & 9 & 34 & 55 & 9.2 & 0.010 & 0.008 \\
53 운송·자재 이동 & 44 & 4 & 16 & 24 & 9.1 & 0.029 & 0.004 \\

\bottomrule
\end{longtable}
\noindent 주: (1) 한국 매칭, (2) global에만 공개, (3) global에서도 미공개. 평균 exposure는 \texttt{observed\_exposure}의 단순 평균. ``--''는 해당 직업이 없음.

\subsection{observed exposure와 같은 원인}

셋째, 매칭되지 않는 직업과 observed exposure가 0인 직업은 대부분 겹친다. 매칭되지 않은 414개 가운데 323개(78\%)가 exposure 0이다. 이는 우연이 아니다. observed exposure는 업무 사용량이 100건 미만인 과업을 0으로 처리하고, 그 사용량은 AEI와 같은 Claude 사용 자료에서 온다(2.2절). 따라서 ``Claude 대화가 거의 없는 직업''이라는 한 가지 사실이 두 자료에 각각 다른 모습으로 나타난다. AEI에서는 셀이 공개되지 않고, observed exposure에서는 값이 0이 된다.

표~\ref{tab:bins}는 이를 exposure 구간별로 보여 준다. exposure가 0인 411개 직업 가운데 한국과 매칭되는 직업은 21.4\%다. 반면 exposure가 0.2를 넘는 직업은 88\% 이상이 매칭된다.

\begin{longtable}{L{4.4cm}rrrrrr}
\caption{observed exposure 구간별 매칭(2026년 5월)}\label{tab:bins}\\
\toprule
observed & 직업 & & & & 한국 매칭 & 한국 \\
exposure 구간 & 수 & (1) & (2) & (3) & 비율(\%) & pct(\%) \\
\midrule
\endfirsthead
\toprule
구간 & 수 & (1) & (2) & (3) & 매칭(\%) & 한국 pct \\
\midrule
\endhead
% 04_analysis_kor_soc_exposure_coverage.do가 생성. 직접 고치지 말 것.
0 & 411 & 88 & 154 & 169 & 21.4 & 5.42 \\
0 초과 0.05 이하 & 94 & 49 & 37 & 8 & 52.1 & 4.21 \\
0.05 초과 0.10 이하 & 68 & 48 & 19 & 1 & 70.6 & 10.69 \\
0.10 초과 0.20 이하 & 64 & 51 & 13 & 0 & 79.7 & 10.08 \\
0.20 초과 0.40 이하 & 84 & 74 & 10 & 0 & 88.1 & 43.68 \\
0.40 초과 & 35 & 32 & 3 & 0 & 91.4 & 20.79 \\

\bottomrule
\end{longtable}

다만 두 자료의 기준이 같지는 않다. exposure가 0이면서 한국과 매칭되는 직업이 88개 있고, 이들의 한국 대화 비중은 5.42\%다. 이 차이는 다음 세 가지로 설명된다. observed exposure는 이론 등급 $\beta < 0.5$인 과업도 0으로 둔다. 또 2025년 8월$\cdot$11월 자료를 쓰므로 2026년 5월과 시점이 다르다. 그리고 AEI \texttt{pct}에는 업무용이 아닌 대화도 포함된다.

% ============================================================
\section{대화량 기준으로 본 손실}
% ============================================================

직업 수로는 절반이 넘게 빠지지만 대화량 기준의 손실은 작다. 2026년 5월 한국에 공개된 세부 직업의 \texttt{pct} 합계는 98.08\%다. 매칭되지 않은 414개 직업은 한국에 공개되지 않은 직업이다. 따라서 이들이 차지하는 한국 대화 비중은 공개되지 않은 나머지 1.92\% 안에 있다. 이 1.92\%에는 \texttt{job\_exposure.csv}에 없는 직업도 포함되므로, 실제 비중은 이보다 작거나 같다.

반대 방향의 손실이 오히려 더 크다. 한국에 공개된 365개 직업 가운데 23개는 \texttt{job\_exposure.csv}에 없다. 이 직업들의 한국 대화 비중은 3.21\%다(표~\ref{tab:reverse}, \ref{tab:kornot}). 결국 매칭된 342개 직업이 한국 대화의 94.87\%를 차지한다.

\begin{longtable}{L{8.0cm}rr}
\caption{AEI 공개 직업 중 \texttt{job\_exposure}에 없는 직업(2026년 5월)}\label{tab:reverse}\\
\toprule
항목 & 한국 & global \\
\midrule
\endfirsthead
\toprule
항목 & 한국 & global \\
\midrule
\endhead
% 04_analysis_kor_soc_exposure_coverage.do가 생성. 직접 고치지 말 것.
공개 SOC 세부 직업(6자리) 수 & 365 & 614 \\
\quad job\_exposure에 있음 & 342 & 578 \\
\quad job\_exposure에 없음 & 23 & 36 \\
job\_exposure에 없는 직업의 pct 합계(\%) & 3.21 & 4.23 \\

\bottomrule
\end{longtable}

\begin{longtable}{lL{8.6cm}rrr}
\caption{한국 공개 직업 중 \texttt{job\_exposure}에 없는 23개 직업(2026년 5월)}\label{tab:kornot}\\
\toprule
 & & 8자리 & 한국 & global \\
코드 & 직업(AEI 명칭) & 코드 수 & pct(\%) & pct(\%) \\
\midrule
\endfirsthead
\toprule
코드 & 직업 & 코드 수 & 한국 pct & global pct \\
\midrule
\endhead
% 04_analysis_kor_soc_exposure_coverage.do가 생성. 직접 고치지 말 것.
25-3031 & Substitute Teachers, Short-Term & 1 & 0.95 & 1.38 \\
29-1212 & Cardiologists & 1 & 0.61 & 0.84 \\
21-1014 & Mental Health Counselors & 1 & 0.38 & 0.43 \\
41-3091 & Sales Representatives of Services, Except Advertising, Insurance, Financial Services, and Travel & 1 & 0.20 & 0.34 \\
13-1023 & Purchasing Agents, Except Wholesale, Retail, and Farm Products & 1 & 0.17 & 0.08 \\
13-1082 & Project Management Specialists & 1 & 0.17 & 0.18 \\
29-2036 & Medical Dosimetrists & 1 & 0.13 & 0.16 \\
29-1214 & Emergency Medicine Physicians & 1 & 0.08 & 0.09 \\
31-1122 & Personal Care Aides & 1 & 0.08 & 0.08 \\
13-2023 & Appraisers and Assessors of Real Estate & 1 & 0.06 & 0.04 \\
25-9042 & Teaching Assistants, Preschool, Elementary, Middle, and Secondary School, Except Special Education & 1 & 0.06 & 0.06 \\
11-9072 & Entertainment and Recreation Managers, Except Gambling & 1 & 0.05 & 0.06 \\
11-2033 & Fundraising Managers & 1 & 0.04 & 0.05 \\
19-4044 & Hydrologic Technicians & 1 & 0.04 & 0.05 \\
53-1042 & First-Line Supervisors of Helpers, Laborers, and Material Movers, Hand & 1 & 0.04 & 0.05 \\
13-1022 & Wholesale and Retail Buyers, Except Farm Products & 1 & 0.03 & 0.08 \\
25-2056 & Special Education Teachers, Elementary School & 1 & 0.03 & 0.05 \\
53-1044 & First-Line Supervisors of Passenger Attendants & 1 & 0.03 & 0.02 \\
39-7011 & Tour Guides and Escorts & 1 & 0.02 & 0.03 \\
25-2055 & Special Education Teachers, Kindergarten & 1 & 0.01 & 0.02 \\
25-9043 & Teaching Assistants, Special Education & 1 & 0.01 & 0.01 \\
39-4012 & Crematory Operators & 1 & 0.01 & 0.01 \\
51-2022 & Electrical and Electronic Equipment Assemblers & 1 & 0.01 & 0.05 \\

\bottomrule
\end{longtable}

이 23개는 모두 SOC 2018 세부 직업 코드이므로 코드 버전 차이 때문에 빠진 것은 아니다. Project Management Specialists(\texttt{13-1082}), Cardiologists(\texttt{29-1212}), Fundraising Managers(\texttt{11-2033})처럼 SOC 2018에서 새로 생겼거나 다시 나뉜 직업이 많다. 이런 직업은 O*NET의 과업 자료나 \textcite{eloundou-2023-gpts-are-gpts-an-early}의 이론 등급이 없어 observed exposure가 계산되지 않았을 가능성이 있다. 다만 이는 직업 목록을 보고 한 추정이며, 원 논문과 자료에서 확인하지 않았다.

% ============================================================
\section{분석 시 유의점}
% ============================================================

\begin{enumerate}[leftmargin=*]
  \item \textbf{미공개는 0이 아니다.} AEI에서 공개되지 않은 직업의 대화량은 ``공개 기준 미만''이다. 이를 0으로 두면 756개 직업을 모두 쓸 수 있다. 하지만 작은 양을 0으로 잘라내게 되고, observed exposure의 0과 같은 원인으로 생긴 0이 양변에 함께 들어간다.
  \item \textbf{결측으로 두면 표본이 선택된다.} 미공개 직업을 결측으로 두면 표본이 342개로 줄어든다. 이 표본은 exposure가 높은 직업 쪽으로 치우친다. 매칭된 직업의 exposure 평균은 0.144이고, 매칭되지 않은 직업은 약 0.022다. 대분류로는 건설$\cdot$생산$\cdot$운송 등 현장 직업이 대부분 빠진다. 고용$\cdot$임금과 결합한 회귀에서는 이 선택이 결과에 영향을 줄 수 있다.
  \item \textbf{두 방식을 함께 보고한다.} 미공개 직업을 0으로 둔 결과와 결측으로 둔 결과를 함께 추정해 견고성을 비교하는 것을 권한다. 미공개 몫(1.92\%)을 미공개 직업에 고르게 나누는 방식도 대안이 될 수 있다.
  \item \textbf{\texttt{job\_exposure}에 없는 23개 직업.} 한국 대화의 3.21\%가 결합에서 빠진다. 같은 broad occupation(6자리의 마지막 자리를 0으로 바꾼 코드)에 속한 직업들의 평균 exposure로 채우는 방법을 검토할 수 있다.
  \item \textbf{단위의 의미.} AEI 직업은 대화 과업의 직업이다. observed exposure 역시 과업 사용에서 만든 직업 지표다. 두 지표를 고용 통계(직업 종사자)와 결합할 때는 ``과업의 직업''과 ``사람의 직업''의 차이를 해석에 반영해야 한다.
\end{enumerate}

% ============================================================
\section{재현}
% ============================================================

Stata 17에서 \texttt{code/04\_analysis\_kor\_soc\_exposure\_coverage.do}를 실행한다. 이 do 파일은 두 원자료(\texttt{data/raw/usage/anthropic\_economic\_index/})를 직접 읽어 본 문서의 표(\texttt{results/table/20\_*.tex})와 엑셀 파일 \texttt{20\_kor\_soc\_exposure\_coverage.xlsx}를 만든다. 엑셀의 \texttt{occupations} 시트에는 756개 직업 각각의 집단과 한국$\cdot$global \texttt{pct}가 들어 있다. 매칭 키는 AEI O*NET-SOC 코드의 앞 7자리와 \texttt{occ\_code}다.

\printbibliography[title={참고문헌}]

\end{document}
